\documentclass[11pt]{article}
\usepackage{amsmath, amssymb}
\usepackage{graphicx}
\usepackage{booktabs}
\usepackage{hyperref}
\usepackage{geometry}
\geometry{margin=1in}

\title{Size Factor Strategies in Differential Expression Analysis:\\
Global vs. Per-Comparison Approaches}
\author{CRISPYx Technical Documentation}
\date{December 2024}

\begin{document}

\maketitle

\section{Overview}

Size factors normalize for differences in sequencing depth across cells. The choice of \textit{which} cells to use when computing size factors significantly affects differential expression results. This document compares two approaches:

\begin{enumerate}
    \item \textbf{Global size factors}: Compute once from all cells, reuse for each comparison
    \item \textbf{Per-comparison size factors}: Compute separately for each perturbation vs. control subset
\end{enumerate}

\section{Mathematical Formulation}

\subsection{DESeq2 Median-of-Ratios Method}

For a count matrix $K$ with cells $i$ and genes $j$, the size factor $s_i$ for cell $i$ is:

\begin{equation}
s_i = \text{median}_j \left( \frac{K_{ij}}{\exp\left( \frac{1}{n} \sum_{i'} \log K_{i'j} \right)} \right)
\end{equation}

where the denominator is the geometric mean of gene $j$ across all cells in the normalization set.

\subsection{Global Size Factors (CRISPYx Default)}

Compute size factors once using \textbf{all $N$ cells} in the dataset:

\begin{equation}
s_i^{\text{global}} = \text{median}_j \left( \frac{K_{ij}}{\exp\left( \frac{1}{N} \sum_{i'=1}^{N} \log K_{i'j} \right)} \right)
\end{equation}

For each perturbation $p$ comparison, use the pre-computed $s_i^{\text{global}}$ for cells in that comparison.

\subsection{Per-Comparison Size Factors (PyDESeq2 Default)}

For each perturbation $p$, compute size factors using only cells in the comparison subset $C_p = \{i : \text{cell } i \in \text{perturbation } p \text{ or control}\}$:

\begin{equation}
s_i^{(p)} = \text{median}_j \left( \frac{K_{ij}}{\exp\left( \frac{1}{|C_p|} \sum_{i' \in C_p} \log K_{i'j} \right)} \right), \quad i \in C_p
\end{equation}

\section{Empirical Comparison}

Testing on the Adamson subset dataset (1716 cells, 100 perturbations):

\begin{table}[h]
\centering
\caption{Size factor statistics by approach}
\begin{tabular}{lcc}
\toprule
\textbf{Metric} & \textbf{Global} & \textbf{Per-Comparison} \\
\midrule
Cells per normalization & 1716 (all) & $\sim$966 (varies) \\
Consistency across comparisons & Identical & Varies \\
Computation cost & $O(1)$ & $O(P)$ where $P$ = perturbations \\
\bottomrule
\end{tabular}
\end{table}

\begin{table}[h]
\centering
\caption{CRISPYx vs PyDESeq2 agreement by size factor approach}
\begin{tabular}{lcc}
\toprule
\textbf{Metric} & \textbf{Global SF} & \textbf{Per-Comparison SF} \\
\midrule
LFC Pearson correlation & 0.95 & 0.99 \\
P-value Pearson correlation & 0.60 & 0.99 \\
Top-100 overlap (LFC) & 94\% & 99\% \\
Top-100 overlap (p-value) & 89\% & 99\% \\
\bottomrule
\end{tabular}
\end{table}

\section{Biological Considerations}

\subsection{CRISPR Screen Context}

In CRISPR screens, perturbations affect specific biological pathways but should not systematically alter global RNA content. This has implications:

\begin{enumerate}
    \item \textbf{Library size is independent of perturbation}: Cells with gene knockouts maintain similar total RNA content
    \item \textbf{Global normalization is appropriate}: Since perturbation doesn't affect sequencing depth, normalizing across all cells captures true library size differences
    \item \textbf{Per-comparison normalization may introduce bias}: Using only perturbed + control cells means size factors are influenced by which genes happen to be affected by the perturbation
\end{enumerate}

\subsection{When Each Approach is Preferred}

\textbf{Global size factors} (CRISPYx default):
\begin{itemize}
    \item CRISPR knockout/knockdown screens
    \item When comparing many perturbations to the same control
    \item When computational efficiency matters
    \item When consistency across comparisons is desired
\end{itemize}

\textbf{Per-comparison size factors}:
\begin{itemize}
    \item When exact PyDESeq2 parity is required for validation
    \item Bulk RNA-seq where each sample is a separate experiment
    \item When perturbations might affect global RNA content (e.g., CRISPRa activation screens)
\end{itemize}

\section{Implementation in CRISPYx}

\subsection{API Usage}

\begin{verbatim}
# Default: Global size factors (efficient, consistent)
nb_glm_test(adata, ..., size_factor_scope="global")

# PyDESeq2 parity: Per-comparison size factors
nb_glm_test(adata, ..., size_factor_scope="per_comparison")
\end{verbatim}

\subsection{Benchmark Configuration}

In \texttt{benchmarking/config/*.yaml}:

\begin{verbatim}
# Enable optional per-comparison size factor method
optional_methods:
  - crispyx_de_nb_glm_sf_per
\end{verbatim}

This runs both:
\begin{itemize}
    \item \texttt{crispyx\_de\_nb\_glm}: Default with global size factors
    \item \texttt{crispyx\_de\_nb\_glm\_sf\_per}: Variant with per-comparison size factors
\end{itemize}

\section{Conclusion}

The choice between global and per-comparison size factors involves a trade-off:

\begin{itemize}
    \item \textbf{Global}: More efficient, more consistent, biologically appropriate for most CRISPR screens
    \item \textbf{Per-comparison}: Matches PyDESeq2 exactly, may be preferred when validating against published results
\end{itemize}

CRISPYx defaults to global size factors for efficiency and biological appropriateness, while providing the \texttt{size\_factor\_scope="per\_comparison"} option for exact PyDESeq2 parity when needed.

\end{document}
